% ============================================================
% bibitems_noche.tex — TT 2026-B039 · 2026-10-14
% Entradas nuevas generadas por el coordinador a partir de las
% respuestas crudas de Crossref, DataCite y la API de arXiv en
% Docs/insumos/fuentes_bib/. Formato APA, como bibitems_nuevos.tex.
% Las entradas de bibitems_nuevos.tex no se repiten aquí.
% Los PDF leídos están en Docs/fuentes/pdf/ (no versionados).
% ============================================================

  % fuente: fuentes_bib/Lee2017CBISDDSM.json · PDF: Lee2017CBISDDSM.pdf
  \bibitem{Lee2017CBISDDSM}
    Lee, R. S., Gimenez, F., Hoogi, A., Miyake, K. K., Gorovoy, M., \& Rubin, D. L. (2017).
    A curated mammography data set for use in computer-aided detection and diagnosis research.
    \textit{Scientific Data}, 4(1), 170177.
    \url{https://doi.org/10.1038/sdata.2017.177}

  % fuente: fuentes_bib/Clark2013TCIA.json (sin PDF; solo para la atribución de TCIA)
  \bibitem{Clark2013TCIA}
    Clark, K., Vendt, B., Smith, K., Freymann, J., Kirby, J., Koppel, P., Moore, S., Phillips, S., Maffitt, D., Pringle, M., Tarbox, L., \& Prior, F. (2013).
    The Cancer Imaging Archive (TCIA): Maintaining and operating a public information repository.
    \textit{Journal of Digital Imaging}, 26(6), 1045--1057.
    \url{https://doi.org/10.1007/s10278-013-9622-7}

  % fuente: fuentes_bib/Zwanenburg2020IBSI.json · PDF: manual de referencia de IBSI (arXiv:1612.07003v11), no el artículo
  \bibitem{Zwanenburg2020IBSI}
    Zwanenburg, A., Valli\`eres, M., Abdalah, M. A., Aerts, H. J. W. L., Andrearczyk, V., Apte, A., Ashrafinia, S., Bakas, S., Beukinga, R. J., Boellaard, R., Bogowicz, M., Boldrini, L., Buvat, I., Cook, G. J. R., Davatzikos, C., Depeursinge, A., Desseroit, M.-C., Dinapoli, N., Dinh, C. V., \ldots\ L\"ock, S. (2020).
    The Image Biomarker Standardization Initiative: Standardized quantitative radiomics for high-throughput image-based phenotyping.
    \textit{Radiology}, 295(2), 328--338.
    \url{https://doi.org/10.1148/radiol.2020191145}

  % fuente: fuentes_bib/Preskill2018NISQ.json · PDF: Preskill2018NISQ.pdf
  \bibitem{Preskill2018NISQ}
    Preskill, J. (2018).
    Quantum computing in the NISQ era and beyond.
    \textit{Quantum}, 2, 79.
    \url{https://doi.org/10.22331/q-2018-08-06-79}

  % fuente: fuentes_bib/McClean2018Barren.json · PDF: McClean2018Barren.pdf (arXiv)
  \bibitem{McClean2018Barren}
    McClean, J. R., Boixo, S., Smelyanskiy, V. N., Babbush, R., \& Neven, H. (2018).
    Barren plateaus in quantum neural network training landscapes.
    \textit{Nature Communications}, 9(1), 4812.
    \url{https://doi.org/10.1038/s41467-018-07090-4}

  % fuente: fuentes_bib/Mari2020Transfer.json · PDF: Mari2020Transfer.pdf (arXiv)
  \bibitem{Mari2020Transfer}
    Mari, A., Bromley, T. R., Izaac, J., Schuld, M., \& Killoran, N. (2020).
    Transfer learning in hybrid classical-quantum neural networks.
    \textit{Quantum}, 4, 340.
    \url{https://doi.org/10.22331/q-2020-10-09-340}

  % fuente: fuentes_bib/SchuldKilloran2022Advantage.json · PDF: SchuldKilloran2022Advantage.pdf (arXiv)
  \bibitem{SchuldKilloran2022Advantage}
    Schuld, M., \& Killoran, N. (2022).
    Is quantum advantage the right goal for quantum machine learning?
    \textit{PRX Quantum}, 3(3), 030101.
    \url{https://doi.org/10.1103/PRXQuantum.3.030101}

  % fuente: fuentes_bib/Shaydulin2022Bandwidth.json · PDF: Shaydulin2022Bandwidth.pdf (arXiv)
  \bibitem{Shaydulin2022Bandwidth}
    Shaydulin, R., \& Wild, S. M. (2022).
    Importance of kernel bandwidth in quantum machine learning.
    \textit{Physical Review A}, 106(4), 042407.
    \url{https://doi.org/10.1103/PhysRevA.106.042407}

  % fuente: fuentes_bib/Canatar2023Bandwidth.json (DataCite) · PDF: Canatar2023Bandwidth.pdf
  % Se cita el preprint de arXiv: la versión publicada (TMLR) no se verificó.
  \bibitem{Canatar2023Bandwidth}
    Canatar, A., Peters, E., Pehlevan, C., Wild, S. M., \& Shaydulin, R. (2022).
    Bandwidth enables generalization in quantum kernel models.
    \textit{arXiv preprint} arXiv:2206.06686.
    \url{https://doi.org/10.48550/arXiv.2206.06686}

  % fuente: fuentes_bib/Kubler2021Inductive.json (DataCite) y journal_ref de arXiv: NeurIPS 34 (2021). Páginas sin verificar.
  \bibitem{Kubler2021Inductive}
    K\"ubler, J. M., Buchholz, S., \& Sch\"olkopf, B. (2021).
    The inductive bias of quantum kernels.
    En \textit{Advances in Neural Information Processing Systems 34 (NeurIPS 2021)}.
    \url{https://doi.org/10.48550/arXiv.2106.03747}

  % fuente: fuentes_bib/Bowles2024Benchmarking_arxiv.xml (API de arXiv) · PDF: Bowles2024Benchmarking.pdf
  \bibitem{Bowles2024Benchmarking}
    Bowles, J., Ahmed, S., \& Schuld, M. (2024).
    Better than classical? The subtle art of benchmarking quantum machine learning models.
    \textit{arXiv preprint} arXiv:2403.07059.
    \url{https://doi.org/10.48550/arXiv.2403.07059}

  % fuente: fuentes_bib/Le2011FRQI.json · PDF: Le2011FRQI.pdf (versión publicada). En línea el 2010-04-17; volumen de 2011.
  \bibitem{Le2011FRQI}
    Le, P. Q., Dong, F., \& Hirota, K. (2011).
    A flexible representation of quantum images for polynomial preparation, image compression, and processing operations.
    \textit{Quantum Information Processing}, 10(1), 63--84.
    \url{https://doi.org/10.1007/s11128-010-0177-y}

  % fuente: fuentes_bib/Zhang2013NEQR.json · PDF: Zhang2013NEQR.pdf (versión publicada)
  \bibitem{Zhang2013NEQR}
    Zhang, Y., Lu, K., Gao, Y., \& Wang, M. (2013).
    NEQR: A novel enhanced quantum representation of digital images.
    \textit{Quantum Information Processing}, 12(8), 2833--2860.
    \url{https://doi.org/10.1007/s11128-013-0567-z}

  % fuente: fuentes_bib/Parigi2025QImR_publicada.json (Crossref) y Parigi2025QImR.json (DataCite, arXiv:2507.22039)
  % PDF: Parigi2026QImR.pdf es la versión de arXiv fechada el 15 de enero de 2026.
  \bibitem{Parigi2026QImR}
    Parigi, M., Khosrojerdi, M., Caruso, F., \& Banchi, L. (2026).
    Analysis of quantum image representations for supervised classification.
    \textit{AVS Quantum Science}, 8(1), 013801.
    \url{https://doi.org/10.1116/5.0296376}


  % fuente: Crossref 10.1056/NEJM199412013312206 (2026-10-16) · texto: Docs/fuentes/txt/Elmore1994/ (T9)
  \bibitem{Elmore1994}
    Elmore, J. G., Wells, C. K., Lee, C. H., Howard, D. H., \& Feinstein, A. R. (1994).
    Variability in radiologists' interpretations of mammograms.
    \textit{New England Journal of Medicine}, 331(22), 1493--1499.
    \url{https://doi.org/10.1056/NEJM199412013312206}

  % fuente: Crossref 10.1016/j.jacr.2021.04.021 (2026-10-16) · texto: Docs/fuentes/txt/Monticciolo2021ACR/ (T9)
  \bibitem{Monticciolo2021ACR}
    Monticciolo, D. L., Malak, S. F., Friedewald, S. M., Eby, P. R., Newell, M. S., Moy, L.,
    Destounis, S., Leung, J. W. T., Hendrick, R. E., \& Smetherman, D. (2021).
    Breast cancer screening recommendations inclusive of all women at average risk: Update from the ACR and Society of Breast Imaging.
    \textit{Journal of the American College of Radiology}, 18(9), 1280--1288.
    \url{https://doi.org/10.1016/j.jacr.2021.04.021}

  % PENDIENTE AL INTEGRAR: AlvarezEstevez2024QuantumKernelBench ya está publicado (Crossref, 2026-10-16):
  % Alvarez-Estevez, D. (2025). Benchmarking quantum machine learning kernel training for classification
  % tasks. IEEE Transactions on Quantum Engineering, 6, 1--15. https://doi.org/10.1109/TQE.2025.3541882
  % Sustituir el \bibitem de Technical_Report.tex (hoy cita el preprint de arXiv de 2024).
